
This study measured the sample efficiency advantage of permutation-equivariant weight-space encoders over plain encoders in a controlled, shared-split experiment on the ModelZooDataset CIFAR-10 CNN zoo. The principal findings are:

\begin{enumerate}
\item GNN-NFN achieves an efficiency ratio of approximately $47\times$ over flat-MLP (conservative lower bound: 6.$8\times )$: it reaches 90\% of its peak $R^2$ at $N\approx 147$ training models while flat-MLP requires the full dataset (~7,000 models) to reach 90\% of its own peak $R^{2}=0.886$.
\end{enumerate}

\begin{enumerate}
\item The efficiency advantage is data-regime dependent. At N=100, PermAug ($R^{2}=0.138)$ outperforms structural equivariance (GNN-NFN $R^{2}=-0.016)$ --- a single-seed finding requiring multi-seed replication. At $N\geq 250$, GNN-NFN dominates. The crossover occurs between N=100 and N=250, suggesting a minimum-data threshold of approximately 150--250 models for equivariant graph encoders on the CIFAR-10 CNN zoo.
\end{enumerate}

\begin{enumerate}
\item Permutation equivariance is verified to floating-point precision for GNN-NFN (max\_diff=1.$80\times 10^{-6}$, 10,000 checks on trained CIFAR-10 zoo models) and DWSNets (max\_diff=7.$45\times 10^{-9}$, 1,000 checks on synthetic MLP weights). The gap ratio versus flat-MLP (max\_diff=5.$59\times 10^{-2})$ is approximately 7.5 million for DWSNets and 31,000 for GNN-NFN.
\end{enumerate}

\begin{enumerate}
\item At full training scale, all encoders converge within 5\% $R^2$ ($gap\leq 0.038)$, consistent with the expressivity equivalence theorem of Dayan et al. [2026].
\end{enumerate}

Three follow-up directions are directly motivated by these results. (1) Multi-seed replication of H-M3 at N=100 with 10 seeds: low cost, addresses the most statistically fragile finding. (2) LR sweep for GNN-NFN at N=100 to distinguish underfitting from learning-rate mismatch as the mechanism for the GNN-NFN N=100 failure. (3) Extension to the MNIST MLP zoo with DWSNets to test generalizability of the efficiency ratio across zoo types and encoder architectures.

The minimum-data threshold identified here --- approximately 150--250 models --- provides practitioners with a concrete criterion for encoder selection under zoo-size constraints. Practitioners with fewer than 150 models should prefer augmentation-based approaches; those with 250 or more models should consider structurally equivariant encoders.

